\section{Experimental Evaluation}

% [proposal] Research questions
We structure the evaluation of Agrasandhani around the following research questions (RQs), which directly reflect the failure modes of naive forwarding and the design goals of smart aggregation under MQTT and degraded networks. \textbf{RQ1 (Bandwidth vs. timeliness):} How much does hybrid batching and latest-per-sensor compaction reduce gateway-to-dashboard bandwidth compared to forwarding every message, and what latency/tail-latency cost does this introduce under normal and bursty loads? \textbf{RQ2 (Reliability-induced duplicates):} Under MQTT QoS~1 (``at least once'' delivery), which may introduce duplicate deliveries via retransmissions \cite{mqtt311}, how effective is message-ID-based deduplication at reducing duplicate-induced overhead while preserving correct state at the dashboard? \textbf{RQ3 (Robustness under impairments):} Under deterministic last-hop impairments (bandwidth caps, loss, delay, jitter, and short outages) injected with Linux traffic control (e.g., \texttt{tc netem}) \cite{jurgelionis2011_netem_empirical}, which gateway variant best maintains dashboard stability as measured by data loss, update jitter, and tail latency? \textbf{RQ4 (User-visible freshness):} How does last-known-good caching with an explicit age indicator and TTL affect the freshness distribution (age of displayed values) during and after impairments, and what trade-offs emerge between reduced communication and increased staleness? \textbf{RQ5 (Ablation and attribution):} Across an ablation ladder (V0--V4), which individual mechanism (batching, compaction, adaptive publish rate, or last-known-good semantics) accounts for the largest improvements for each metric?

\subsection{Evaluation Setting}
% [proposal] Evaluation plan
We will evaluate Agrasandhani using an end-to-end local testbed consisting of (i) a Python/asyncio sensor simulator publishing JSON messages \texttt{\{sensor\_id, msg\_id, ts\_sent, value, type\}} to an MQTT broker, (ii) a smart gateway subscribing to the raw topic and publishing aggregated updates to an output topic, and (iii) a lightweight dashboard backend that subscribes to aggregated updates and streams them to a browser UI over WebSockets. The primary comparison is an ablation ladder of gateway variants: V0 (naive forward-every-message), V1 (hybrid batching with flush every $T$ ms or $K$ messages), V2 (V1 plus two-stage deduplication and latest-per-sensor compaction), V3 (V2 plus adaptive publish rate by adjusting $T$ based on publish health signals such as send time and queue backlog), and V4 (V3 plus last-known-good dashboard semantics with explicit age and TTL). For each variant, we will run MQTT under QoS~0 and QoS~1 to expose reliability/duplicate effects; MQTT v3.1.1 defines QoS~0 as ``at most once'' and QoS~1 as ``at least once'' delivery \cite{mqtt311}. We will sweep workload intensity by varying the number of simulated sensors and their rate schedule (including deterministic burst phases), while using a mixed sensor model (slowly changing versus spiky values) to reflect heterogeneous update dynamics. To stress the user-visible last hop, we inject deterministic, repeatable degradations primarily on the gateway-to-dashboard path using the application-layer impairment proxy, with optional host-level Linux traffic control (bandwidth caps plus \texttt{tc netem} for loss, delay, jitter, and outage-like disruptions) in the same last-hop context \cite{jurgelionis2011_netem_empirical}. This placement matches the measured downstream evidence, which is collected from proxy-level frame and byte counters. Each condition will be executed for multiple trials (e.g., 2--3) with fixed random seeds to support repeatability. We will report end-to-end latency from \texttt{ts\_sent} to dashboard receipt/display (including p50/p95/p99), message rate (msgs/sec) at each stage, bandwidth usage (bytes/sec) on the gateway egress, and data loss defined via missing/late updates inferred from per-sensor \texttt{msg\_id} sequences. In addition, we will measure user-visible stability via update jitter (inter-arrival variance) and freshness via the age distribution of displayed values (with TTL-based staleness marking), enabling a rigorous characterization of the trade-offs between communication reduction, timeliness, and dashboard robustness under degraded networks.

\subsection{Evaluation Results}
All quantitative claims in this section are traceable to generated artifacts under \texttt{report/assets} and copied paper assets under \texttt{research\_paper/figures} and \texttt{research\_paper/tables}. The source runs are the final April evidence runs and their companion sweeps listed in \texttt{report/assets/evidence\_manifest.json}.

\begin{figure*}[t]
	\centering
	\includegraphics[width=\linewidth]{figures/main_outage_frame_rate.png}
	\caption{Downstream Frame Rate Over Time During Outage: V0 vs V2 vs V4 (Intel, outage\_5s, QoS0). V0 shows bursty/high downstream frame rate, while V2 and V4 remain much lower and steadier before, during, and after outage. In this evidence set, smart variants reduce downstream frame cadence by roughly 95--97\% versus naive forwarding, stabilizing the visualization stream around outage and recovery.}
	\label{fig:main-outage-frame-rate}
\end{figure*}

\begin{table}[t]
\centering
\caption{Primary Intel summary (QoS~0, means across three replicated trials).}
\label{tab:main-summary}
\scriptsize
\setlength{\tabcolsep}{3pt}
\renewcommand{\arraystretch}{1.03}
\begin{adjustbox}{max width=\columnwidth}
\begin{tabular}{@{}lrr>{\raggedleft\arraybackslash}p{0.16\columnwidth}>{\raggedright\arraybackslash}p{0.26\columnwidth}@{}}
\toprule
Variant & Frames & Bytes & \makecell[r]{Latency p95\\(ms)} & Scenario \\
\midrule
V0 & 132 & 12069 & 31.9 & clean \\
V2 & 5 & 13051 & 276.2 & clean \\
V4 & 6 & 17641 & 607.7 & clean \\
V0 & 132 & 12069 & 121.5 & \nolinkurl{bandwidth_200kbps} \\
V2 & 5 & 13051 & 271.5 & \nolinkurl{bandwidth_200kbps} \\
V4 & 6 & 17641 & 573.7 & \nolinkurl{bandwidth_200kbps} \\
V0 & 128 & 11725 & 115.7 & \nolinkurl{loss_2pct} \\
V2 & 4 & 12129 & 267.0 & \nolinkurl{loss_2pct} \\
V4 & 5 & 16845 & 546.0 & \nolinkurl{loss_2pct} \\
V0 & 116 & 10605 & 119.8 & \nolinkurl{outage_5s} \\
V2 & 4 & 11401 & 279.1 & \nolinkurl{outage_5s} \\
V4 & 6 & 17641 & 541.7 & \nolinkurl{outage_5s} \\
\bottomrule
\end{tabular}
\end{adjustbox}
\end{table}


Figure~\ref{fig:main-outage-frame-rate} is the primary paper figure. It shows the user-visible effect of aggregation on rendering cadence under outage, where baseline V0 remains bursty and smart variants maintain a low-rate stream. The latency CDF for clean QoS0 (V0/V2/V4) is shown in \texttt{research\_paper/figures/intel\_clean\_qos0\_latency\_cdf.png}, and confirms the expected latency-versus-stability tradeoff.

Bandwidth-over-time and message-rate-over-time outage traces are provided in \texttt{research\_paper/figures/intel\_outage\_qos1\_bandwidth\_over\_time.png} and \texttt{research\_paper/figures/intel\_outage\_qos1\_message\_rate\_over\_time.png}. The batch-window tradeoff is shown in \texttt{research\_paper/figures/intel\_v2\_batch\_window\_tradeoff.png}, where larger windows increase latency and reduce frame cadence, with byte volume changing across the sweep.

The V1 vs V2 isolation result appears in \texttt{research\_paper/figures/intel\_v1\_vs\_v2\_isolation.png} and \texttt{research\_paper/tables/intel\_v1\_vs\_v2\_isolation.md}. It shows small or mixed effects beyond batching for this configuration. The adaptive V2 vs V3 impairment comparison in \texttt{research\_paper/figures/intel\_v2\_vs\_v3\_adaptive\_impairment.png} shows a null result under tested thresholds.

Freshness/retention behavior during outage is shown by the age-over-time trace \texttt{research\_paper/figures/intel\_outage\_qos0\_v0\_vs\_v4\_age\_over\_time.png} and table \texttt{research\_paper/tables/intel\_outage\_qos0\_v0\_vs\_v4\_freshness.md}. In this setup, last-known-good improves visibility/retention but not age freshness. QoS0 vs QoS1 side-by-side outcomes are shown in \texttt{research\_paper/figures/intel\_qos\_comparison.png} and \texttt{research\_paper/tables/intel\_qos\_comparison.md}, with minimal differences in this local environment.